\documentclass{article}
\usepackage{graphicx} % Required for inserting 
\usepackage{amssymb}
\usepackage[margin=1in]{geometry}


\title{Projet, MTH8414}
\author{Antoine Habib, Isaac Valin}
\date{\today}

\begin{document}
\maketitle

\section*{Question 1 - using HiGHs}

\subsection*{1.1}
%Décrire les décisions 
\begin{itemize}
    \item variable de décision qui prend en compte deux décision à la fois:
        \subitem Décider si bus $b$ recharge à la période $t$ 
        \subitem Décider la puissance de recharge de bus $b$ à la période $t$
\end{itemize}

\subsection*{1.2}
\paragraph{Esembles}
\begin{itemize}
    \item B - Nombre de bus dans le problème, $b = 1, \dots, B$
    \item T - Nombre de périodes du problème, $t = 0,\dots 48$
\end{itemize}

\paragraph{Paramètres}
\begin{itemize}
    \item $a_{bt}$ - si bus b est en route pendant période $t$, $a_{bt}\in [0,1], \quad \forall b \in B, \quad \forall t \in T$
    \item $e_{bt}$ - energie de bus $b$ à la fin de la période $t$, $e_{bt} \in \mathbb{R}$
    \item $c_{bt}$ - consommation d'énergie de bus $b$ pendant la période $t$, $c_{bt}$ $\in \mathbb {R}, \quad \forall b \in B, \quad \forall t \in T$
    \item $p_t$ -  prix d'énergie de recharge pendant periode $t$. 
    
\end{itemize}


\paragraph{Variables}
$x_{bt}$, la quantité d'énergie de recharge de bus $b$ à la période $t$, $\qquad x_{bt}\in [0,150]$

\paragraph{Contraintes}
\begin{itemize}
    \item $x_{ij} \leq M(1-a_{bt})$, pour $M \geq 150 (kW), \quad \forall b \in B, \quad \forall t \in T$, , Si bus $b$ à période $t$ est sur la route, ne peut pas charger. Si le bus recharge, il ne peut pas être sur la route. 
    \item $x_{bt} \leq 150, , \quad \forall b \in B, \quad \forall t \in T$. Contrainte de charge maximale
    \item $ 90 \leq e_{bt} \leq 450, , \quad \forall b \in B, \quad \forall t \in T$. Contrainte de capcité de batterie et de reserve minimale. 
    \item $e_{b0}=e_{b48}=450$, Contrainte de batterie pleine en début et fin de temps.
    \item $e_{bt+1} = e_{bt} - c_{bt} + x_{bt}, , \quad \forall b \in B, \quad \forall t \in T$. Contrainte de consommation et de rechargement d'énergie. 
    \item 
\end{itemize}

\paragraph{Fonction Objectif}
minimiser les coûts de recharge pour tous les bus sur toutes les périodes : 
$$\min \sum_{b} x_{bt}p_t$$




\end{document}
